% §III Task and the Failure (~0.75 page, Fig. 2). 0729-pure per Q2: no June depletion sweep /
% gate-oracle in body. Failure = naive 0/3 overpress (canonical: EVIDENCE §3 08-11 블록). The
% policy is phenomenologically described here; mechanism diagnosis deferred to §V (probes).
% 08-11 수정: 용어 naive policy로 통일; abort는 external-layer/censored 서술 (§I과 정합);
% 성공 카운트 미표기 ("every evaluation trial"). 로봇 명시 08-11 사용자 결정: Dexmate Vega-1p
% (하드웨어 벤더 명시는 double-anonymous 허용 관행; 필요시 camera-ready 전 재검토).
\section{Task, Setting, and the Failure}
\label{sec:task}

\textbf{Task and platform.}
We study suction case picking on a Dexmate Vega-1p, a dual-arm mobile manipulator: the arm
descends onto a stack of rigid cases, must stop when the suction cup meets the top case, hold
until vacuum seal is confirmed, and lift (Fig.~\ref{fig:setup}). The cases carry battery
cells, which tolerate little compression---over-pressing risks the product, not just the
pick---so deployment enforces a $15$\,N contact-force limit. The margin is tight in both
directions: too little press and the seal fails; a few millimeters too much and contact force
climbs past the limit, since the stack is stiff (${\sim}5$\,N/mm, Sec.~\ref{sec:discussion}). The policy
controls a 7-D end-effector action ($3$-D position and $3$-D rotation increments plus a
suction command) at $15$\,Hz, executed in chunks of $5$ actions (${\sim}330$\,ms between
observations).

\textbf{Observations.}
The policy observes a single head-mounted camera (RGB and colorized depth; no wrist camera) and a
15-D state: end-effector position and quaternion, suction command, vacuum-seal bit, and the 6-D
wrist wrench. The setting is a deliberately clean diagnostic: at head-camera range, case-stack
heights differing by millimeters---the entire success margin---are visually near-indistinguishable
(Fig.~\ref{fig:setup}, right), while the wrench and seal bit signal the contact transition
directly. Every condition the task depends on is \emph{in the state the policy receives}; the
question is whether the policy binds its behavior to it.

\textbf{Demonstrations.}
We collect $100$ training episodes ($19{,}445$ frames) and hold out $6$ validation episodes
($1{,}169$ frames) with a scripted expert designed to \emph{decorrelate} force from its habitual
correlates: each episode descends with suction on, presses to a randomized force target drawn
from $[8,15]$\,N, retreats $5$--$10$\,mm from the contact plane, then hovers and lifts
immediately on seal. Episodes are collected at two different stack heights, so the contact
plane itself moves across the dataset. Every episode thus demonstrates the counterfactual the task requires:
force rises $\rightarrow$ stop $\rightarrow$ move back up, \emph{before} any seal event, at a
force level that varies independently of depth. This is the data-side prescription whose
rationale we develop in Sec.~\ref{sec:discussion}; here it matters that the naive policy and
all conditioned variants are trained on exactly these episodes.

\textbf{The failure.}
Deployed on the robot, the naive policy---a SmolVLA~\cite{smolvla} fine-tuned on these
demonstrations, with the full wrench and seal bit in its state---fails every evaluation trial
the same way: it descends, makes contact, and keeps pressing---never initiating a stop of its
own---until the external safety layer aborts the run ($15.4$--$18.7$\,N measured at abort).
The transition that the demonstrations exhibit in every single episode, driven by the one signal
that observably announces it, is the transition the policy does not perform. The rollout says
the policy is bad; it does not say \emph{why}. The policy might ignore force entirely; it might
respond to force but too weakly; it might respond in a form that collapses precisely when force
exceeds the demonstrated range. These hypotheses have opposite prescriptions (more data?
different sensors? different architecture?), and none of them is distinguishable from training
curves or rollouts. Distinguishing them requires intervening on the signal itself---the
instrument we build next.
